\chapter{HANDOVER Supplementary Case Study}
\label{app:handover}

\section{Objective and Relationship to the Research Questions}
HANDOVER tests whether paired video and force--torque signals can identify the recorded
outcome of a completed human--robot handover. The single-input predictors address
representation choice (SQ1), and averaging their probabilities tests the contribution of
the second input on the same trials (SQ2). For SQ3, controlled blur, sensor noise, and
single-stream fallback show how the baselines behave when input quality or availability
changes. For SQ4, class-specific decisions at a fixed threshold complement the ranking
scores. Missing participant identifiers limit the uncertainty analysis and prevent
participant-disjoint fitting of the proposed learned fusion model.\ev{HANDOVER study
interpretation}{\EvidenceRef{E084}{}}

The HANDOVER dataset paper compared video, force/torque and gripper inputs for
completed-trial failure detection \parencite{Thoduka2024Handover}. Its five-class I3D model
achieved 69.1\% accuracy with video and force/torque, falling to 67.9\% when the gripper
input was added. The authors also investigated temporal action modelling with auxiliary
action supervision. The binary experiment here evaluates a different target and uses AP,
which precludes a numerical comparison with those five-class accuracies.

\section{Dataset and Common-Support Cohort}
The source describes successful and failed robot handovers \parencite{Thoduka2024Handover}.
The evaluation uses the pinned Zenodo v1 release \parencite{Thoduka2024HandoverData}, with
589 native trials involving Toyota HSR and Kinova Gen3 and 17 participants reported by
its datasheet. Native outcomes are success, no approach, outcome~2, and drop. Outcome~2
means no grasp for robot-to-human handover or no release for human-to-robot handover.
The binary target maps success to zero and every nonzero native outcome to failure.
Trial IDs restart across official splits, so identities are split-qualified; an ID is
not a participant surrogate.

The native official split contains 337 training, 101 validation and 151 test trials. The
common-support cohort retains 587: 335 train, 101 validation and all 151 test trials,
after two training trials fail the timing gate. The test has 41 successes and 110
failures, giving positive prevalence 0.728477. Released outcome counts govern the
analysis where a small datasheet discrepancy changes the distribution of outcome~2
versus drop; it does not change the binary failure total. The official participant-based
split is the publisher's claim, not an independently verified identity
separation.\ev{HANDOVER native and admitted cohort}{\EvidenceRef{E080}{}}

\section{Preparation and Input Construction}
The initial full-video analysis excluded 92 test trials because their final frames
extended beyond wrench coverage. Before scoring the models, the observation window was
defined as the intersection of camera and wrench coverage for every trial, class and
robot. The 16 source-clock bin-centre frames and numerical statistics use that same
window. All official test trials could then be assessed. The completed-trial label can
still depend on an event after the end of the shared window.

For the authors' resampling procedure, matching the nearest raw wrench timestamp
reproduced the supplied arrays under a maximum permitted gap of 0.1~s. Video
presentation timestamps and sensor time refer to separate clocks. Six statistics for
each of six force/torque channels give 36 numerical features. The HSR measures wrench
directly, whereas the Gen3 estimates it from joint torques. Fourteen traces met the
timing criterion but were zero on every channel. These trials were included with their
original labels.\ev{HANDOVER common support and wrench semantics}{\EvidenceRef{E081}{}}

\section{Models and Fitting Procedure}
The pretrained visual condition (P) used the frozen ResNet18-vd CNN backbone inside the
RT-DETR-R18vd checkpoint. Spatial and temporal means produced one 896-dimensional vector
per trial. A newly fitted 896--128--1 head trained for ten fixed epochs. Averaging the
frames removed their order. The numerical branch fitted training-only scaling and
balanced L2 logistic regression to the 36 statistics. A fixed HGB control, a
training-prior predictor, and equal averaging of the visual and logistic predictions
completed the five baselines. Each was evaluated under three fixed seeds. Learned
stacking and the quality gates required participant identifiers for their training
predictions.

HANDOVER uses statistical wrench descriptors in place of the REASSEMBLE raw-sequence
temporal branch. Its own data determine the visual head and sensor scaler. None of the
fitted source-classifier components or F2 coefficients were transferred. The planned R
condition required a REASSEMBLE-trained CNN checkpoint that met the HANDOVER selection
criteria. No eligible checkpoint was established for that condition.\ev{HANDOVER model
boundary}{\EvidenceRef{E083}{}}

\section{Evaluation and Results}
The primary metric is failure-positive AP on the same 151 official test trials, computed
separately for each seed. Table~\ref{tab:handover-clean} presents the three-seed mean
and descriptive spread. The deterministic numerical branches repeat identical
predictions across seed labels; their zero spread does not imply independent
replication.

\begin{table}[htbp]
\caption{HANDOVER Test Performance on the Common Set of Trials}\label{tab:handover-clean}
\centering\small

\begin{tabular}{lrr}\toprule Model&Mean AP&Seed SD\\\midrule
Training-prior predictor&0.728477&0.000000\\
Visual&0.771647&0.006411\\
Sensor logistic&0.823624&0.000000\\
Sensor HGB&0.906319&0.000000\\
Averaging&0.803102&0.013458\\\bottomrule
\end{tabular}
\displaynote{Failure is the positive class. All models use the same 151 trials.}
\end{table}
\ev{HANDOVER test comparison}{\EvidenceRef{E082}{}}

Averaging increased AP by 0.031455 relative to vision and reduced it by 0.020522
relative to logistic sensing. All three seed differences against sensing were negative.
The failure-class F1 score also reflects the high proportion of failures. With a threshold of 0.5
the prior predictor calls every trial a failure, giving an F1 score of 0.842912 and zero
specificity for success. Logistic sensing had higher AP than vision but a success
specificity of only 0.170732.

Under the strongest blur, averaging yielded AP of 0.779066, compared with 0.823624 for
intact logistic sensing. At the highest sensor-noise level, averaging yielded 0.796955,
noisy logistic sensing 0.824952 and HGB 0.731344. Removing either input entirely
triggered a fixed fallback to the surviving branch.\ev{HANDOVER stress
boundaries}{\EvidenceRef{E084}{}}

\section{Participant Identity and the Learned-Fusion Restriction}
F2, F2-aug, QG-clean and QG-aug required out-of-fold branch probabilities from
participant-disjoint splits. The mapping from trial to participant was unavailable.
Trial IDs, robot IDs and task labels do not identify individual participants. The
official validation set retained its original role, and none of these four fusion
methods was fitted.

The missing participant mapping also prevents estimation of a participant-cluster
interval. A trial bootstrap would treat repeated observations from the same person as
independent. The variation across seeds describes optimisation on a fixed set of test
trials. Participant independence is reported for the official split by the dataset
authors, but could not be checked from trial-level identities.\ev{HANDOVER grouping
restriction}{\EvidenceRef{E082}{}}

\section{Interpretation of the HANDOVER Results}
On the 151 handovers, logistic sensing ranked failures better than the visual model or
simple averaging. The HGB sensor model achieved the highest clean AP among the five
baselines. These results concern predictions made from the common camera and wrench
observation window.

Whether a learned fusion rule could improve on sensing remains open for this cohort. The
proposed comparison requires a participant mapping or equivalent grouping information to
obtain independent training predictions for the fusion model. Without that comparison
the study provides no estimate of an effect from learned reliability weighting or
transfer of source-trained CNN weights.
